\chapter{Planificación académica general}
\label{chap:planificacion-academica}

\section{Criterio de planificación}

La planificación conserva el calendario académico comprendido entre marzo y diciembre de 2026, pero diferencia dos niveles. El primero reúne los hitos formales de la asignatura: propuesta, entregas del 25 \%, 50 \% y 75 \%, exposición y defensa. El segundo organiza la evolución técnica del producto mediante checkpoints, dependencias entre repositorios y evidencia de cierre.

Esta separación evita medir el avance únicamente por cantidad de código. Cada bloque se considera cerrado cuando cuenta con implementación o documentación, criterio de aceptación, evidencia y conexión con el roadmap general.

\section{Planificación por entregables}

\begingroup
\scriptsize
\setlength{\tabcolsep}{3pt}
\renewcommand{\arraystretch}{1.18}
\begin{longtable}{|p{0.28\textwidth}|p{0.25\textwidth}|p{0.39\textwidth}|}
\caption{Roadmap actualizado para la entrega del 50 \%.}\label{tab:plan-roadmap-50}\\
\hline
\textbf{Entregable} & \textbf{Estado} & \textbf{Evidencia} \\
\hline
\endfirsthead
\hline
\textbf{Entregable} & \textbf{Estado} & \textbf{Evidencia} \\
\hline
\endhead
Fundamentación y entrega del 25 \% & terminado & capítulos 1–9 y propuesta aprobada. \\
\hline
User research inicial & realizado; corte consolidado para la entrega del 50 \% & cinco entrevistas y encuesta cuantitativa con 167 respuestas. \\
\hline
Modelo sagital & entrenado y evaluado; supera el umbral de calidad del proyecto & checkpoint, manifest y métricas sobre partición retenida por paciente. \\
\hline
Modelo axial & entrenado y evaluado; el artifact permanece por debajo del umbral de calidad & checkpoint, manifest y métricas; ver Capítulo 13. \\
\hline
P9 integración multiplanar & terminado & ejecución, persistencia, workspace y reapertura. \\
\hline
P10 seguridad y observabilidad & consolidado para P10.9 & criterios de seguridad del checkpoint P10.9 consolidados; hardening productivo posterior \\
\hline
P10.5-A contrato volumétrico & terminado & commits coordinados en AI Module, backend y frontend. \\
\hline
P10.5-B a P10.5-F evolución del visor volumétrico & implementado & catálogo de series, serving de cortes, shell de visor, vinculación por corte y reapertura E2E Implementados y congelados por checkpoint en sus ramas correspondientes; navegación volumétrica disponible además en la rama principal del frontend. \\
\hline
Documentación del 50 \% & en integración & documento maestro, capítulos revisados, matrices y diagramas. \\
\hline
Google Cloud & despliegue parcial DEV y arquitectura objetivo documentada & Artifact Registry/WIF, Backend Cloud Run DEV y Frontend Cloud Run DEV desplegados; AI privado, Cloud SQL, Secret Manager, storage productivo, DEMO environment y monitoring live pendientes. \\
\hline
Validación profesional cualitativa sobre producto & realizada con alcance acotado & dos sesiones sobre versión funcional/intermedia; observaciones y ajustes priorizados; no equivale a validación clínica ni cruzada. \\
\hline
Entrega del 75 \% y defensa & futura & producto consolidado, resultados, discusión y presentación. \\
\hline
P10.6 clasificador subarticular (Axial T2) & modelo entrenado y evaluado; capacidad automática bloqueada & checkpoint congelado sobre RSNA / LumbarDISC; falta ROI axial automática validada \\
\hline
P10.7 hallazgos degenerativos asistivos por nivel & implementado pero bloqueado para inferencia productiva automatica & checkpoint y contratos existen; \texttt{blocked\_preprocessing\_parity} y \texttt{automaticInferenceEnabled=false}; no se presenta como productivo. \\
\hline
P10.8 preflight y guardas & cerrado sin entrenamiento adicional & verificación de factibilidad y guardas de contrato Cerrado sin entrenamiento adicional; estado congelado en su rama de checkpoint. \\
\hline
P10.9 consolidación de producto (Backend + AI) & congelado para traspaso de Frontend & commits coordinados; Backend\ensuremath{\leftrightarrow}AI validado; contrato público sanitizado; fixtures P10.9 Backend + AI congelados para handoff de Frontend en rama de checkpoint identificable, con commit y evidencia del runner técnico; base posteriormente integrada en Frontend. \\
\hline
Integración P10.9 en Frontend & implementada & contratos P10.9 consumidos por el Frontend, catálogo full-series, visor navegable y reapertura integrada; Patient queda parcial porque el historial/longitudinal visible usa principalmente \texttt{subjectRef}. \\
\hline
Recorrido Full Product UI E2E sobre stack real & validado funcionalmente & Playwright final 11/11 PASS sobre el recorrido demo/producto; no equivale a validacion clinica ni a Patient browser-level completo. \\
\hline
\end{longtable}
\endgroup

\section{Gestión de cambios}

Los cambios de alcance deben registrarse con su origen. La incorporación axial se vincula con user research y factibilidad de datos. El visor volumétrico se vincula con la constatación de que los MHA ya contienen múltiples cortes y con la necesidad de aproximar la experiencia a un lector real. La comparación longitudinal fue implementada y validada funcionalmente en Demo Mode y flujo basado en \texttt{subjectRef}; para entradas reales no-demo depende de una localización anatómica resuelta.

Cada cambio debe actualizar requerimientos, roadmap, arquitectura, contratos, pruebas y documentación. No corresponde modificar solo la interfaz cuando la capacidad depende de nuevos assets o persistencia.

\section{Riesgos de planificación}

Los principales riesgos son: sobrealcance clínico, apertura de funcionalidades antes de cerrar dependencias, falta de evidencia, divergencia entre repositorios, costos cloud no controlados, ausencia de validación cruzada profesional final e inconsistencias entre documentación y producto.

Las mitigaciones son: checkpoints compartidos, commits coordinados, criterios de aceptación, separación entre implementado y planificado, control presupuestario, pruebas por capas y uso del documento maestro como fuente editorial común.

\section{Conclusión de la planificación}

La planificación no reinicia el proyecto después del 25 \%: toma como base los modelos, la integración multiplanar y la arquitectura ya desarrollados. El camino crítico ya no pasa por generar y persistir las vistas previas de todos los cortes ni por integrar los contratos P10.9 en el Frontend, capacidades actualmente implementadas. El tramo pendiente ya no consiste en cerrar el recorrido E2E funcional ni en obtener observaciones cualitativas iniciales sobre el producto, sino en ejecutar la validación cruzada profesional final, completar el análisis económico/negocio, terminar el cierre discursivo de la tesis y avanzar desde el despliegue parcial DEV hacia los recursos Cloud no desplegados. La región de interés axial automática constituye una línea paralela que no condiciona ninguno de esos pasos.
